\documentclass[14pt]{extarticle}
\usepackage{fontspec}
\usepackage{polyglossia}
\usepackage{amsmath}
\usepackage{amssymb}
\usepackage{xcolor}
\usepackage{fancyhdr}
\usepackage{graphicx}
\usepackage{listings}
\usepackage[most]{tcolorbox}
% \usepackage{setspace}
\usepackage{pifont}
\usepackage{enumitem}
\usepackage{titlesec}
\usepackage[bottom]{footmisc}
\usepackage{titling}
\usepackage{minted}
\usepackage{etoolbox}
\usepackage{array}
\usepackage{longtable} % page-breakable tables (must load before polyglossia/bidi RTL setup)
\usepackage{extsizes}
\usepackage{float}
\usepackage{booktabs}
\usepackage{hyperref}
\usepackage[onehalfspacing]{setspace}
\newif\ifwhiteblack
\ifdefined\setwhiteblack
  \setwhiteblack
\fi

\ifwhiteblack
  \usepackage[a4paper, landscape, top=1.25cm, bottom=1.25cm, left=1.25cm, right=1.25cm, headheight=1.5cm, headsep=0.5cm]{geometry}
\else
  \usepackage[a4paper, top=2.7cm, bottom=2.7cm, left=2.5cm, right=2.5cm, headheight=1.5cm, headsep=0.5cm]{geometry}
\fi
\usepackage{tikz}
\usepackage{datetime2}

\usetikzlibrary{automata, positioning, arrows.meta, arrows}
\providecommand{\LTcaptype}{table}

\newfontfamily\emoji{Segoe UI Emoji}

\pagestyle{fancy}

\setmainlanguage[numerals=western]{arabic}
\setotherlanguage{english}
\newfontfamily\arabicfont[Script=Arabic]{Calibri}
\newfontfamily\arabicfonttt[Script=Arabic]{Courier New}



\ifwhiteblack
\lstset{
  language=[Sharp]C,
  numbers=left,
  stepnumber=1,
  numbersep=8pt,
  frame=single,
  basicstyle=\ttfamily\small,
  keywordstyle=\bfseries,
  stringstyle=\ttfamily,
  commentstyle=\itshape
}
\else
\lstset{
  language=[Sharp]C,
  numbers=left,
  stepnumber=1,
  numbersep=8pt,
  frame=single,
  basicstyle=\ttfamily\small,
  keywordstyle=\color{blue},
  stringstyle=\color{red},
  commentstyle=\color{green!50!black}
}
\fi

\newif\ifdetailed
\ifdefined\setdetailed
  \setdetailed
\fi

\newif\ifwithsols
\ifdefined\setwithsols
  \setwithsols
\fi

\ifwhiteblack
  \AtBeginDocument{
    \LARGE
    \usemintedstyle{bw}
  }
\fi

% Default minted look: compact code that wraps long lines.
% Not applied in black-and-white landscape mode, which already uses \LARGE.
\ifwhiteblack\else
  \setminted{fontsize=\small, breaklines=true}
\fi

% unified tcolorboxes for articles
\ifwhiteblack
\tcbset{colback=white, colframe=black, fonttitle=\bfseries, boxrule=1pt}
\newtcolorbox{boxDef}[1][]{title={{\emoji📘} تعريف\ifx\\#1\\\else ~#1\fi :}}
\newtcolorbox{boxExercise}[1][]{title={{\emoji🧩} تمرين\ifx\\#1\\\else ~#1\fi :}}
\newtcolorbox{boxExample}[1][]{title={{\emoji📝} مثال\ifx\\#1\\\else ~#1\fi :}}
\newtcolorbox{boxNote}[1][]{title={{\emoji✨} ملاحظة\ifx\\#1\\\else ~#1\fi :}}
\newtcolorbox{boxAttention}[1][]{title={{\emoji🔔} تنبيه\ifx\\#1\\\else ~#1\fi :}}
\newtcolorbox{boxWarning}[1][]{title={{\emoji⚡} ملاحظة هامة\ifx\\#1\\\else ~#1\fi :}}
\newtcolorbox{boxSolution}[1][]{title={{\emoji✅} حل\ifx\\#1\\\else ~#1\fi :}}
\newtcolorbox{boxSymbol}[1][]{title={{\emoji🔣} رمز\ifx\\#1\\\else ~#1\fi :}}
\newtcolorbox{boxHint}[1][]{title={{\emoji💡} تلميح\ifx\\#1\\\else ~#1\fi :}}
\else
\tcbset{colback=white, colframe=black, fonttitle=\bfseries, boxrule=0.8pt}
\newtcolorbox{boxDef}[1][]{colback=blue!5!white,colframe=blue!75!black,
  title={{\emoji📘} تعريف\ifx\\#1\\\else ~#1\fi :}}
\newtcolorbox{boxExercise}[1][]{colback=cyan!5!white,colframe=cyan!70!black,
  title={{\emoji🧩} تمرين\ifx\\#1\\\else ~#1\fi :}}
\newtcolorbox{boxExample}[1][]{colback=yellow!5!white,colframe=orange!90!black,
  title={{\emoji📝} مثال\ifx\\#1\\\else ~#1\fi :}}
\newtcolorbox{boxNote}[1][]{colback=gray!10!white,colframe=black,
  title={{\emoji✨} ملاحظة\ifx\\#1\\\else ~#1\fi :}}
\newtcolorbox{boxAttention}[1][]{colback=magenta!10!white,colframe=magenta!80!black,
  title={{\emoji🔔} تنبيه\ifx\\#1\\\else ~#1\fi :}}
\newtcolorbox{boxWarning}[1][]{colback=red!5!white,colframe=red!75!black,
  title={{\emoji⚡} ملاحظة هامة\ifx\\#1\\\else ~#1\fi :}}
\newtcolorbox{boxSolution}[1][]{colback=green!5!white,colframe=green!60!black,
  title={{\emoji✅} حل\ifx\\#1\\\else ~#1\fi :}}
\newtcolorbox{boxSymbol}[1][]{colback=purple!5!white,colframe=purple!70!black,
  title={{\emoji🔣} رمز\ifx\\#1\\\else ~#1\fi :}}
\newtcolorbox{boxHint}[1][]{colback=teal!5!white,colframe=teal!60!black,
  title={{\emoji💡} تلميح\ifx\\#1\\\else ~#1\fi :}}
\fi


\ifwhiteblack
\tcbset{simplecode/.style={ colback=white, colframe=black, boxrule=1pt, arc=2pt, left=4pt,right=4pt,top=4pt,bottom=4pt}}
\else
\tcbset{simplecode/.style={ colback=gray!5, colframe=black!50, boxrule=0.4pt, arc=2pt, left=4pt,right=4pt,top=4pt,bottom=4pt}}
\fi
% Compact code boxes: tighter line spacing (the preamble uses onehalfspacing) and small vertical gaps.
% Must come AFTER the simplecode definition above.
\tcbset{simplecode/.append style={fontupper=\setstretch{0.95}, top=4pt, bottom=4pt, before skip=6pt, after skip=6pt}}
\newenvironment{boxCode}{\begin{tcolorbox}[simplecode]}{\end{tcolorbox}}

% Code cell used inside tables (keeps code left-to-right inside RTL tables).
% Use inside a fixed-width p{...} column: \begin{pycell}\begin{minted}...\end{minted}\end{pycell}
\newtcolorbox{pycell}{frame empty,colback=white,boxsep=0pt,left=2pt,right=2pt,top=2pt,bottom=2pt,width=\linewidth}

\newcolumntype{C}[1]{>{\centering\arraybackslash}p{#1}}

\makeatletter
\DTMnewdatestyle{dotted}{%
  \renewcommand{\DTMdisplaydate}[4]{\two@digits{##3}.\two@digits{##2}.##1}%
  \renewcommand{\DTMDisplaydate}{\DTMdisplaydate}%
}
\makeatother
\DTMsetdatestyle{dotted}

% redefine spaces after titles
\makeatletter
\renewcommand{\@maketitle}{%
  \begin{center}
    {\huge \bfseries \@title \par}%
    \vskip 0.2em % space between title and author
    {\large \@author \par}%
    % \vskip 0.2em % space between author and date
    % {\normalsize \@date \par}%
    {\small \hfill \begin{english}\DTMtoday\end{english} \par}%
  \end{center}
}
\makeatother

\fancyhf{} % clear default
\fancypagestyle{plain}{
  \fancyhf{}
  \fancyhead[L]{مدرسة التسامح الشاملة}
  % \fancyhead[L]{\includegraphics[height=1cm]{../../../images/logoTasamoh.png}}
  \fancyhead[R]{الأستاذ محمود اغبارية}
  \fancyfoot[C]{\thepage}
}

\fancyhead[L]{مدرسة التسامح الشاملة}
\fancyhead[R]{الأستاذ محمود اغبارية}
\fancyfoot[C]{\thepage}

\setcounter{tocdepth}{3} % only section subsection and subsubsection in TOC

\newcommand{\importpdfpage}[2]{%
    \thispagestyle{empty}
    \newgeometry{left=0cm,right=0cm,top=0cm,bottom=0cm}
    \includegraphics[page=#2,width=0.95\paperwidth,height=0.95\paperheight,keepaspectratio]{#1}
    \restoregeometry
}

\newcommand{\insertFullImg}[1]{%
    \noindent
    \makebox[\textwidth][c]{
        \includegraphics[width=0.9\paperwidth,keepaspectratio]{#1}%
    }%
}%


\title{وظيفة بيتية 4: المتغيرات ومكتبة \textenglish{Math}}

\begin{document}

\maketitle
\thispagestyle{fancy}

\begin{enumerate}[itemsep=3em]

    \item
    اكتب برنامجًا يستقبل من المستخدم عددًا من 4 منازل (مثلاً 1234).
    البرنامج يقوم ببناء عدد جديد بحيث يتم التبديل بين منزلة العشرات ومنزلة المئات (1324).
    البرنامج يطبع العدد الجديد، ثم يطبع الفارق المطلق بين العددين.
    \ifdetailed
    \begin{boxExample}
    \begin{english}
    \begin{minted}{csharp}
Please enter a 4-digit number:
1234
New number = 1324
Absolute difference = 90
    \end{minted}
    \end{english}
    \end{boxExample}
    \ifwithsols
    \begin{boxSolution}
    \begin{english}
    \begin{minted}{csharp}
int num = int.Parse(Console.ReadLine());

int units = num % 10;
int tens = (num / 10) % 10;
int hundreds = (num / 100) % 10;
int thousands = num / 1000;

int newNum = thousands * 1000 + tens * 100 + hundreds * 10 + units;
int diff = Math.Abs(num - newNum);

Console.WriteLine("New number = " + newNum);
Console.WriteLine("Absolute difference = " + diff);
    \end{minted}
    \end{english}
    \end{boxSolution}
    \clearpage
    \fi
    \fi

    \item
    اكتب برنامجًا يستقبل عددًا يمثل إجمالي عدد الأيام. البرنامج يقوم بتحويل هذا العدد إلى سنوات (اعتبر السنة 365 يومًا)، وأشهر (اعتبر الشهر 30 يومًا)، وأيام متبقية.
    \ifdetailed
    \begin{boxExample}
    \begin{english}
    \begin{minted}{csharp}
Please enter total days:
800
Years = 2
Months = 2
Days = 10
    \end{minted}
    \end{english}
    \end{boxExample}
    \ifwithsols
    \begin{boxSolution}
    \begin{english}
    \begin{minted}{csharp}
int totalDays = int.Parse(Console.ReadLine());

int years = totalDays / 365;
int remainingDays = totalDays % 365;

int months = remainingDays / 30;
int days = remainingDays % 30;

Console.WriteLine("Years = " + years);
Console.WriteLine("Months = " + months);
Console.WriteLine("Days = " + days);
    \end{minted}
    \end{english}
    \end{boxSolution}
    \fi
    \fi
    \clearpage

    \item
    اكتب برنامجًا يستقبل وقت بداية حدث ما كعدد صحيح من 4 منازل بنظام 24 ساعة (مثلاً 1430 تعني 14:30)، ويستقبل مدة الحدث بالدقائق (مثلاً 135 دقيقة).
    البرنامج يحسب وقت انتهاء الحدث ويطبعه بنفس الصيغة (HHMM).
    \ifdetailed
    \begin{boxExample}
    \begin{english}
    \begin{minted}{csharp}
Please enter start time (HHMM):
1430
Please enter duration in minutes:
135
End time (HHMM) = 1645
    \end{minted}
    \end{english}
    \end{boxExample}
    \ifwithsols
    \begin{boxSolution}
    \begin{english}
    \begin{minted}{csharp}
int startTime = int.Parse(Console.ReadLine());
int duration = int.Parse(Console.ReadLine());

int startHours = startTime / 100;
int startMins = startTime % 100;

int totalStartMins = startHours * 60 + startMins;
int totalEndMins = totalStartMins + duration;

int endHours = totalEndMins / 60;
int endMins = totalEndMins % 60;

int endTime = endHours * 100 + endMins;
Console.WriteLine("End time (HHMM) = " + endTime);
    \end{minted}
    \end{english}
    \end{boxSolution}
    \clearpage
    \fi
    \fi

    \item
    اكتب برنامجًا يستقبل من المستخدم عددًا عشريًّا يمثّل قيمة $x$.
    البرنامج يحسب ويطبع قيمة الدالة التالية عند النقطة $x$ (البرنامج يقرب النتيجة إلى فاصلتين):
    $$ f(x) = \sqrt{|x^3 - 1|} $$
    \ifdetailed
    \begin{boxExample}
    \begin{english}
    \begin{minted}{csharp}
Please enter x:
3
f(x) = 5.1
    \end{minted}
    \end{english}
    \end{boxExample}
    \ifwithsols
    \begin{boxSolution}
    \begin{english}
    \begin{minted}{csharp}
double x = double.Parse(Console.ReadLine());

double cube = Math.Pow(x, 3);
double absolute = Math.Abs(cube - 1);
double result = Math.Sqrt(absolute);

double resultTimes100 = result * 100;
double resultRounded = Math.Round(resultTimes100) / 100.0;

Console.WriteLine("f(x) = " + resultRounded);
    \end{minted}
    \end{english}
    \end{boxSolution}
    \fi
    \fi

\end{enumerate}

\end{document}
